\section{Matrix reconstruction and the two obstruction quotients}
\label{v7:app:matrix-reconstruction}

The new arguments use algebraic correspondences and finite matrix
identities in distinct ways. We record the identities so that the
computations have a precise scope. They check the displayed linear
algebra, not the existence of the missing cycles.

\subsection{The two-stage obstruction}

For every smooth projective \(X\), the inclusions
\(\Alg^p(X)\subseteq\Ab^p(X)\subseteq\Hdg^p(X)\) give
the exact sequence
\begin{equation}\label{v7:eq:two-stage-exact}
0\longrightarrow\frac{\Ab^p(X)}{\Alg^p(X)}
 \longrightarrow D_{\Alg}^p(X)
 \longrightarrow D_{\Ab}^p(X)\longrightarrow0.
\end{equation}
The first map is induced by inclusion, so is injective. The second
sends the coset of \(\alpha\) modulo \(\Alg\) to its coset modulo
\(\Ab\); it is surjective, and its kernel consists exactly of
classes in \(\Ab\) modulo \(\Alg\). This proves exactness without
choosing a splitting. All three maps commute with algebraic
correspondences.

Every class in the left term is a finite sum of images of classes in
\(D_{\Alg}(A)\), where \(A\) is abelian: use the defining expression
\(\sum\Gamma_{i*}\beta_i\) for \(\Ab\), and pass to the quotient.
Replacing \(\beta_i\) by an algebraic class changes its image by an
algebraic class, so these quotient maps are well defined. Universal
\Fii{} makes every such source quotient zero; it therefore kills
the left term. Universal \Fiii{} kills the right term. Their
conjunction kills the middle term, which is the rational Hodge
conjecture. This is a description of the two obligations, not a
proof that either quotient vanishes.

\begin{figure}[ht]
\centering
\begin{tikzpicture}[>=Stealth,node distance=6mm,
 box/.style={draw,rounded corners,align=center,inner sep=6pt,font=\small}]
\node[box,fill=blue!5] (a) {Abelian-source defect\\\(\Ab^p/\Alg^p\)};
\node[box,fill=gray!6,right=of a] (b) {Algebraicity defect\\\(\Hdg^p/\Alg^p\)};
\node[box,fill=orange!7,right=of b] (c) {Transfer defect\\\(\Hdg^p/\Ab^p\)};
\draw[->] (a)--node[above]{\(\subseteq\)}(b);
\draw[->>] (b)--(c);
\node[below=7mm of a,font=\small] {universal F2 kills this term};
\node[below=7mm of c,font=\small] {F3\( {}'\) kills this term};
\end{tikzpicture}
\caption{The exact sequence \eqref{v7:eq:two-stage-exact}.
The middle obstruction vanishes when both outer terms vanish.
The diagram does not assert that they already vanish.}
\label{v7:fig:two-defects}
\end{figure}

\subsection{The projector matrix}

Write \(v_i=hD_i\), and let
\(\ell_j(\alpha)=\int_X\alpha D_jh^{d-3}\). Then
\(\ell_j(v_i)=Q_{ji}\). If \(V\) is the column matrix with
columns \(v_i\) and \(L\) has rows \(\ell_j\), then
\[
 LV=Q,\qquad P=VQ^{-1}L,\qquad
 P^2=VQ^{-1}(LV)Q^{-1}L=P.
\]
Moreover \(PV=V\), \(L(1-P)=0\), and \(P(1-P)=0\).
These identities give the coefficient extraction and the complementary
projection in Proposition~\ref{v7:prop:primitive-hodge-projector}.
The proof of that proposition identifies the complement with primitive
classes on \(\Hdg^2(X)\).\footnote{The finite identities alone
would only identify \(\ker L\). The hard Lefschetz argument on
Hodge classes is what identifies that kernel with the primitive
Hodge subspace. No claim about the full non-Hodge cohomology follows
from the finite matrix test.}

For the completely numerical model
\[
 Q=\begin{pmatrix}4&1\\1&-2\end{pmatrix},\qquad
 Q^{-1}=\frac19\begin{pmatrix}2&1\\1&-4\end{pmatrix},
\]
take \(V\) to be the first two coordinate vectors of \(\QQ^4\)
and \(L=(Q\mid 0)\). Then \(P=\operatorname{diag}(1,1,0,0)\).
The last two coordinates model a remaining quotient; the example
does not identify them with classes on an actual variety.

\subsection{Triangular recovery and traces}

The projective-bundle coefficient equations have upper triangular
matrix \(T\) with entries \(T_{jk}=u_{k-j}\) for \(k\ge j\),
\(T_{jj}=1\), and zero for \(k<j\). If \(N=T-I\), then
\(N^{r+1}=0\), and the exact inverse is
\[
 T^{-1}=I-N+N^2-\cdots+(-1)^rN^r.
\]
Thus the recursive coefficient formulas involve only finite sums,
compositions and multiplication by the stated algebraic classes.

For a generically finite map, put \(A=f^*\), \(B=f_*\) on
the cohomology or on either defect quotient. The identity \(BA=dI\)
gives \(e=AB/d\), \(e^2=e\), \(eA=A\), and
\(B(1-e)=0\). Hence the trace-zero complement cannot absorb a
nonzero class pulled back from the base. The finite checks reproduce
these identities in rectangular rational matrices satisfying
\(BA=dI\). They do not verify the geometric degree identity, which
is proved using the projection formula in the text.

The supplemental program \texttt{checks/check\_transfer.py} uses
exact rational arithmetic for all three calculations. Its reported
cases concern those identities only. The new geometric arguments
and the universal conjecture have not been formalized in Lean.

\subsection{Coordinate bounds and denominators}

\begin{lemma}
\label{v7:new:lem:cramer}
Let \(\gamma\in N_0^{-1}\ZZ^b\), with
\(\|\gamma\|_\infty\le G\), \(G\ge1\). Suppose \(\gamma\) belongs
to the rational span of integral vectors \(v_i\) satisfying
\(\|v_i\|_\infty\le M\), where \(M\ge1\) is an integer.
It can be written using at most \(b\) such vectors, with coefficients
\[
 |q_i|\le b!\,G M^{b-1},
\]
whose denominators all divide
\[
 N_0\operatorname{lcm}(1,2,\ldots,b!M^b).
\]
\end{lemma}

\begin{proof}
Choose \(r\le b\) independent columns spanning \(\gamma\), and
choose \(r\) rows on which their matrix \(A\) is invertible.
Its determinant is a nonzero integer and
\(|\det A|\le r!M^r\le b!M^b\).
Cramer's rule gives
\(q_i=\det A_i/\det A\), with
\(|\det A_i|\le r!G M^{r-1}\le b!G M^{b-1}\).
Also \(N_0\det A_i\) is integral. Taking the indicated least
common multiple clears all possible determinants.
\end{proof}

On a relatively compact analytically trivialized part of \(S\),
degree bounds on effective cycles do give coordinate bounds of this
kind. Choose smooth closed forms dual to a flat integral basis and
a smooth relative Kahler form representing \(c_1(L)\). Their norms
are uniformly bounded on the compact part. Integrating over a
complex subvariety bounds each class coordinate by a fixed constant
times its \(L\)-degree. This gives a concrete lattice explanation
for automatic denominators. It supplies no bound on the degrees
needed to represent the missing mixed class.
